%!TEX root = ../book.tex
% Ukrainian translation of Correll et al., Introduction to Autonomous Robots
% Source: github.com/Correll-Lab/Introduction-to-Autonomous-Robots
% File: chapters/sensors.tex
% CC BY-NC-SA license per repo.
% Translation by auxiliary local run (Anthropic), 2026-06-01.
\chapter{Сенсори}\label{chap:sensors}

Роботи — це системи, що сприймають, діють і обчислюють. Досі ми вивчали базові фізичні принципи руху, тобто локомоції та маніпуляції. Тепер потрібно зрозуміти основні принципи сенсорів робота, що надають необхідні дані для прийняття рішень і керування.

Цілі цього розділу:
\begin{itemize}
\item подати огляд сенсорів, поширених у роботизованих системах;
\item викласти фізичні принципи їх роботи;
\item з'ясувати механізми, відповідальні за невизначеність у сенсорному висновуванні.
\end{itemize}

Історично розвиток роботизованих сенсорів спричинений не лише робототехнікою, а й суміжними галузями: транспортом (автомобільний, морський, авіація), пристроями безпеки промисловості, сервоприводами в радіокерованих іграшках, а останнім часом — смартфонами, віртуальною реальністю та ігровими консолями. Ці галузі переважно зробили ``екзотичні'' сенсори доступними за низькою ціною завдяки масовим застосуванням. Наприклад, акселерометри та гіроскопи зараз широко використовуються у смартфонах і коштують менше долара; XBox зробила 3D-сенсорику глибини (Kinect) доступною; сучасні легкові автомобілі містять велику кількість сенсорів без значного збільшення вартості.

\begin{mdframed}
Подумайте про сенсори, з якими ви взаємодієте щодня. Які сенсори у вашому телефоні, на кухні, в автомобілі?
\end{mdframed}

Сенсори важко класифікувати лише за областю застосування. Більшість задач виграє від кожного можливого джерела інформації. Наприклад, локалізації можна досягти, вимірюючи кількість обертів колеса \emph{енкодером} (\ref{sec:sensors:encoders}), але оцінювання стає точнішим з додаванням акселерометрів (\cref{sec:sensors:inertia}) або візуальних сенсорів (\cref{chap:vision}). Усі підходи відрізняються за точністю і типом наданих даних, але жоден з них не розв'язує задачу локалізації самостійно.

\section{Термінологія}\label{sec:sensors:terminology}

При роботі з сенсорами важливо точно визначати такі терміни, як ``швидкість'' і ``роздільна здатність'', а також додаткову таксономію.

Розрізняють \emph{активні} та \emph{пасивні} сенсори. Активні сенсори\index{активний сенсор} випромінюють певну енергію та вимірюють реакцію середовища. Пасивні сенсори\index{пасивний сенсор} вимірюють енергію, що надходить із середовища. Більшість дистанційних сенсорів (окрім стерео-зору) — активні. Акселерометр, компас, кнопка-натискач — пасивні.

\emph{Діапазон}\index{діапазон (сенсора)} — \emph{різниця} між верхньою і нижньою межею вимірюваної величини. Відрізняється від \emph{динамічного діапазону}\index{динамічний діапазон (сенсора)} — \emph{відношення} між найвищим і найнижчим значеннями (зазвичай у логарифмічному масштабі — ``децибел''\index{децибел}). Мінімальна відстань між двома значеннями, які сенсор може розрізнити — \emph{роздільна здатність}\index{роздільна здатність (сенсора)}.

\emph{Точність}\index{точність (сенсора)} (accuracy) — різниця між (середнім) виходом $m$ і справжньою величиною $v$:
\begin{equation}
\text{accuracy} = 1 - \frac{|m-v|}{v}.
\end{equation}
\emph{Прецизійність}\index{прецизійність (сенсора)} (precision) — відношення діапазону до статистичної дисперсії сигналу. Прецизійність — це \emph{повторюваність}, а точність описує \emph{систематичну похибку} фізики сенсора.
\begin{figure}
    \centering
    \def\svgwidth{0.9\textwidth}
    \import{./figs/}{precisionvsaccuracy.pdf_tex}
    \caption{Хрестик — справжнє значення. Зліва направо: ані прецизійно, ані точно; прецизійно, але не точно; точно, але не прецизійно; точно і прецизійно. \label{fig:precision}}
\end{figure}
Наприклад, GPS зазвичай прецизійний у межах кількох метрів, але точний лише в межах десятків. Це особливо проявляється при зміні конфігурації супутників.

Швидкість надання вимірювань — \emph{пропускна здатність} сенсора\index{пропускна здатність (сенсора)}. Сенсор з пропускною здатністю 10~Гц дає сигнал десять разів на секунду.

\subsection{Пропріоцепція vs екстероцепція}\label{sec:sensors:proprioception}

\emph{Пропріоцепція}\index{пропріоцепція} — сприйняття внутрішнього стану робота (кути суглобів, швидкості, внутрішні моменти й сили).

\emph{Екстероцепція}\index{екстероцепція} — сприйняття всього зовнішнього: стану світу, невизначеності та правильна дія. Останнім часом увагу приділяють \emph{проксимальним} сенсорам — для вимірювання середовища безпосередньо біля тіла робота, а не лише \emph{дистальним} (камери, акустичні).

\section{Сенсори положення суглобів}\label{sec:sensors:encoders}

Найважливіший пропріоцептивний сенсор — \emph{енкодер}\index{енкодер}. Енкодери використовуються для виміру положення і швидкості суглоба, а також сили (у поєднанні з пружиною). Поділяються на інкрементальні (відносні; переважно в мобільних роботах) і абсолютні (переважно в маніпуляторах). Загалом покладаються на магнітний або оптичний маяк, що обертається з мотором, і сенсор, що рахує проходи. Найпоширеніший в робототехніці — \emph{квадратурний енкодер}\index{квадратурний енкодер}, оптичний, що використовує патерн, який обертається з мотором, і оптичний сенсор для реєстрації переходів чорний/білий (\cref{fig:encoders}).

\begin{figure}
    \centering
    \includegraphics[width=0.3\textwidth]{figs/encoderdisk.png}
    \fontsize{7pt}{9pt}\selectfont
    \def\svgwidth{0.3\textwidth}
    \import{./figs/}{quadraturencoder.pdf_tex}
    \includegraphics[width=0.3\textwidth]{figs/absoluteencoder.png}
    \caption{Зліва направо: патерн квадратурного енкодера; вихідний сигнал (рух вперед); патерн абсолютного енкодера (код Грея).}
    \label{fig:encoders}
\end{figure}

Один сенсор достатній для виявлення положення і швидкості, але не дозволяє визначити напрямок руху. Квадратурні енкодери мають два сенсори (A, B), що реєструють чергуючий патерн з відстанню у чверть фази. Якщо A випереджає B, диск обертається за годинниковою стрілкою; якщо B випереджає A, проти. Абсолютні енкодери використовують код Грея\index{код Грея} — патерн, у якому між сусідніми сегментами змінюється лише один біт.

\section{Сенсори его-руху}\label{sec:sensors:inertia}

Вимірювання конфігурації суглобів обмежене статичними спостереженнями. Воно не дозволяє роботу виявити, чи він зараз рухається або прискорюється — особливо важливо для динамічно стійких роботів (гуманоїдів, квадрокоптерів). Рух можна оцінити через принцип \emph{інерції}\index{інерція}.

\subsection{Акселерометри}

\emph{Акселерометр}\index{акселерометр} можна уявити як масу на демпфованій пружині. За закону Гука $F=kx$ вимірюємо силу, яку маса спричиняє на пружину. З $F=ma$ обчислюємо прискорення $a$ маси $m$. На Землі це приблизно $9{,}81\,\mathrm{м/с^2}$.
На практиці ці системи реалізують мікроелектромеханічними пристроями (MEMS) — наприклад, консольною балкою, відхилення якої вимірюється ємнісним сенсором. Акселерометри вимірюють до трьох осей трансляційних прискорень. Виведення абсолютного положення вимагає подвійного інтегрування, що вносить значний шум — це робить оцінку положення лише з акселерометра непрактичною.
Однак гравітація дає сталий вектор прискорення, тому акселерометри добре оцінюють орієнтацію відносно гравітації (\emph{крен} і \emph{тангаж}).

\subsection{Гіроскопи}\label{sec:gyroscopes}

\emph{Гіроскоп}\index{гіроскоп} — електромеханічний пристрій, що вимірює кутову швидкість і (у деяких конфігураціях) орієнтацію. Комплементарний акселерометру. Класично — це диск, що вільно обертається в системі осей; рух системи зберігається інерційним моментом. Дискові гіроскопи досі використовуються, але важко мініатюризуються.

Варіація — \emph{гіроскоп кутової швидкості} (rate gyro)\index{rate gyro}, що вимірює кутову/обертальну швидкість.
У \emph{оптичному} гіроскопі лазерний промінь розщеплюють надвоє і пускають по круговому шляху в протилежних напрямках. При обертанні системи проти напрямку одного з променів цей промінь долатиме трохи довшу відстань — це дає вимірюваний фазовий зсув на приймачі, пропорційний \emph{швидкості обертання}.
Малорозмірні оптичні rate gyros непрактичні, але MEMS rate gyros широко доступні — використовують іншу технологію: масу, підвішену на пружинах. Маса активно вібрує, тому при обертанні сенсора зазнає сили Коріоліса, пропорційної бічній швидкості й кутовій швидкості.

Rate-гіроскопи вимірюють кутову швидкість навколо трьох осей, інтегрування дає абсолютну орієнтацію. Акселерометр (3 осі трансляції) + гіроскоп (3 осі обертання) → інформація про рух у всіх шести ступенях свободи. Разом з магнітометром (компасом), що дає абсолютну орієнтацію за курсом (yaw), це формує \emph{інерціальний вимірювальний модуль}\index{інерціальний вимірювальний модуль} (IMU\index{IMU}). Ця комбінація особливо потужна: акселерометр + гіроскоп дають комплементарну інформацію про крен/тангаж, магнітометр + гіроскоп — про курс. Це лежить в основі систем визначення кутової орієнтації (AHRS\index{AHRS}) через сенсорне злиття.

\section{Вимірювання сили}
Вимір фізичних сил взаємодії — критично важливий для робототехніки. Дає змогу делікатно піднімати ягоду, безпечно взаємодіяти з людиною через дотик.

При поєднанні мотора, енкодера й пружини отримуємо \emph{серійний пружний актуатор} (Series Elastic Actuator)~\cite{pratt1995series}\index{серійний пружний актуатор} — обертальні/лінійні енкодери можуть служити простими сенсорами сили чи моменту через закон Гука. Інший метод — вимірювання струму на кожному суглобі: знаючи позу, обчислюємо очікувані сили й струми, а похибки відповідають додатковим силам.

\begin{figure}
    \centering
    \def\svgwidth{0.9\textwidth}
    \import{./figs/}{ftsensor.pdf_tex}
    \caption{Сенсор сили/моменту (F/T) передає силу і момент між двома ланками робота через три металеві стрижні, що з'єднують внутрішню втулку із зовнішнім кільцем. Кожен стрижень обладнаний тензодатчиком з кожного боку — 12 сенсорів усього. \label{fig:ftsensor}}
\end{figure}

Найточніший і найпоширеніший — \emph{сенсор сили/моменту} (F/T)\index{сенсор сили/моменту}. Механічний пристрій, що детектує одну або більше компонент шестивимірного \emph{wrench}-у (3D-сили і 3D-моменту). Більшість комерційних F/T-сенсорів використовують \emph{тензодатчики}\index{тензодатчик} — металеву (провідну) фольгу, що змінює форму під дією wrench і змінює електричний опір. Типова конфігурація: внутрішня втулка з суцільного металу підвішена у зовнішньому кільці через три симетричних прямокутних металевих стрижні, кожен з 4 тензодатчиками — разом 6 сигналів (попарно), з яких обчислюємо сили й моменти в 3D.

Обмеження F/T: 1) висока вартість через високу точність виготовлення; 2) розмір (типово розмір зап'ястя людини); 3) низьке відношення сигнал/шум; 4) низька пропускна здатність/реактивність; 5) одинична точка даних, розріджена в просторі і часі (особливо для маніпулятора з багатьма точками контакту).

\subsection{Вимірювання тиску або дотику}
Щоб частково усунути ці обмеження, робототехніки працюють над вимірюванням тиску, прикладеного до поверхні робота.

Людський дотик — найстарший, найбільший і один із найважливіших органів чуття.

\emph{Сенсор тиску} здатний детектувати або контакт/зіткнення як бінарну величину (сенсор дотику), або градієнт тиску. Вертикальний тиск пропорційний 1D-силі по нормалі до сенсора. Сенсори тиску переважно ємнісні (як тачскрін смартфона): два провідні шари, розділені ізолятором; при тиску відстань зменшується, ємність змінюється.

У порівнянні з F/T, сенсори тиску дають менше інформації (1D vs 6D), але мають: 1) високу реактивність; 2) високу щільність сенсингу (до десятків сенсорів на см$^2$); 3) низьку вартість; 4) легкість мініатюризації.

Людський дотик не обмежується лише тиском — також високочастотною інформацією (вібраціями). Роботи можуть відтворювати це інтегруванням акселерометрів або мікрофонів у м'який перетворювач та класифікацією спектральної інформації.

В ідеалі — \emph{штучна шкіра}\index{штучна шкіра}, що поєднує різні модальності (тиск, текстура, температура, світло) — поки не отримала широкого застосування.

\section{Сенсори вимірювання відстані}
Існує плавний перехід від пропріоцептивних до екстероцептивних сенсорів — вимір внутрішнього стану тісно пов'язаний із зовнішнім середовищем при контакті. Для дистанційного дослідження середовища критично виміряти відстань до окремих об'єктів.

Малий форм-фактор і низька ціна світлочутливих напівпровідників призвели до поширення оптичних сенсорів на різних фізичних ефектах: відбиття, фазовий зсув, час прольоту. Інші принципи — радіо (``радар'') і звук.

\subsection{Відбиття}
Найпростіший принцип: чим ближче об'єкт, тим більше світла/радіо/звуку він відбиває. Для незалежності від кольору об'єкта зазвичай використовують \emph{інфрачервоне} світло; звук залежить від поверхневих властивостей.

Сенсор має дві компоненти: випромінювач і приймач (вимірює силу відбитого сигналу). Типова відповідь IR-сенсора показана на рис.~\ref{fig:epuckir} — насичена на малих відстанях, квадратично зменшується далі.

\begin{figure}
    \centering
    \includegraphics[width=0.8\textwidth]{figs/epuckirsensor.png}
    \caption{Реальна відповідь IR-сенсора відстані як функція відстані. Одиниці навмисно безрозмірні.}
    \label{fig:epuckir}
\end{figure}

\subsection{Фазовий зсув}\label{sec:phaseshiftsensors}
На відстанях більших за короткі відбиттєве вимірювання неточне. Лазерні сенсори вимірюють \emph{фазовий зсув} відбитої хвилі. Випромінене світло модулюється довжиною хвилі, що перевищує максимальну вимірювану дистанцію.

Сучасні \emph{лазерні сканери}\index{лазерний далекомір} випромінюють на 5~МГц. Зі швидкістю світла $3\cdot 10^5$~км/с довжина хвилі — 60~м, що дає сканеру корисний діапазон до 30~м.
Знаючи форму випроміненої хвилі, обчислюємо фазовий зсув, звідки — відстань.

Лазерні сканери в поєднанні з обертальними дзеркалами утворюють \emph{лазерні скануючі далекоміри} (\emph{LiDAR}\index{LiDAR}). Системи з 64 скануючими лазерами широко використовуються в автономному водінні.

\subsection{Час прольоту (time-of-flight)}
Найточніше вимірювання — за часом прольоту світла. Світло рухається швидко ($3\cdot 10^8$~м/с), тому потрібна високошвидкісна електроніка, що вимірює періоди менше наносекунди для сантиметрової точності.
На практиці поєднують приймач з дуже швидким електронним затвором, що працює на тій же частоті, що й випромінений імпульс. Знаючи час, виводимо відстань за кількістю фотонів, що повернулись.

\subsubsection{Ультразвукові сенсори відстані}\label{sec:sensors:sound}
Вимір часу прольоту значно простіший для звукових хвиль (~344~м/с у повітрі). Ультразвуковий сенсор випромінює імпульс і вимірює час до повернення.
Через нижчу швидкість звуку — компроміс між дальністю і пропускною здатністю: довша дальність вимагає довшого очікування.

Перевага ультразвукових перед оптичними: замість променя — конус з кутом розкриття $20$--$40^\circ$. Можуть детектувати малі перешкоди без необхідності прямого попадання променя. Це робить їх сенсором вибору в авто-парковочних системах сучасних автомобілів.

\section{Сенсори глобальної пози}\label{sec:sensors:globalpose}
Досі ми обговорювали сенсори, що вимірюють положення суглобів, кутову швидкість, трансляційне прискорення, сили взаємодії та відстань до об'єктів відносно власної пози. Для надійної навігації роботи потребують поняття \emph{світової системи координат}.

Локалізація об'єкта через тріангуляцію існувала ще в давні часи — моряки орієнтувалися за зорями. Найдосконаліша система — \emph{Глобальна система позиціонування} (GPS)\index{GPS}. Супутники на орбіті оснащені знанням свого точного положення і синхронізованими годинниками; вони транслюють радіосигнали, закодовані часом випромінювання. Приймачі обчислюють відстань до кожного супутника, порівнюючи час випромінювання й час прийому. Оскільки невідоме як положення $(x,y,z)$, так і часова різниця між годинниками — потрібно 4 супутники для отримання ``фіксу''. Початковий фікс — кілька хвилин, потім — кілька разів на секунду. GPS-вимірювання недостатньо прецизійні й точні для робототехнічних застосувань і вимагають злиття з іншими сенсорами (наприклад, IMU).

Існують також \emph{закриті GPS-рішення} — активні або пасивні маяки з відомими локаціями. Пасивні (інфрачервоні відбивні наклейки, 2D-штрихкоди) детектуються камерами і дають позу з відомих розмірів. Активні випромінюють радіо, ультразвук або їх комбінацію — для оцінки відстані. У цій області особливо поширене ультраширокосмугове радіо (UWB).

\section*{Висновки на винос}
\begin{itemize}
\item Більшість сенсорів робота вирішує задачу визначення пози робота або локалізації/розпізнавання об'єктів у його оточенні.
\item Кожен сенсор має переваги й недоліки, що квантифікуються в його діапазоні, прецизійності, точності й пропускній здатності. Тому стійкі рішення досягаються лише через поєднання сенсорів з різними принципами роботи.
\item Твердотільні сенсори (без механічних частин) можна мініатюризувати й дешево виробляти. Це уможливило появу доступних IMU і 3D-сенсорів глибини — основи локалізації та розпізнавання об'єктів у масових роботизованих системах.
\end{itemize}

\section*{Вправи}\small
\begin{enumerate}
\item Лазерний сканер з кутовою роздільною здатністю 0.01~рад і максимальною дальністю 5.6~м. Яка мінімальна відстань $d$ від робота до об'єкта шириною 1~см, щоб гарантовано вловити його хоча б одним променем? Апроксимуйте відстань між променями довжиною дуги.
\item Чому пропускна здатність ультразвукового сенсора різко падає при збільшенні динамічного діапазону, а лазерного — ні?
\item Ви проектуєте автономний електромобіль для перевезення вантажів. Ви турбуєтесь про вартість, обираючи між лазерним сканером і ультразвуковим. Дальність — 15~м. Ультразвук швидкість 300~м/с.
\begin{enumerate}
\item Час від ультразвука до перешкоди 15~м.
\item Час від лазера.
\item Якщо рухаєтесь до перешкоди — який сенсор дасть ближче до реального вимірювання й чому?
\end{enumerate}
\item Виберіть навчальну робот-платформу і складіть список її сенсорів.
\item Сконструюйте простий range scanner з ультразвука на сервомоторі. Чи бачите прості ознаки — кути, отвори?
\item Знайдіть DIY robotics магазини в інтернеті. Які сенсори вони пропонують? Які інтерфейси?
\item Виберіть фізичний сенсор. Чи можете спроектувати експеримент для характеризації його прецизійності й точності?
\item Спроектуйте сенсор для виявлення вільного простору в посилці для e-commerce.
\item Спроектуйте автономний візок, що автоматично пристиковується до полиць у середовищі.
\item Ви проектуєте контролер для гри ``Ratslife''. Яку інформацію надає середовище і який сенсор потрібен?
\end{enumerate}\normalsize
